\documentclass[10pt]{amsart}
\usepackage[a4paper,margin=23mm]{geometry}
\usepackage{amsmath,amssymb,amsthm,mathtools,hyperref,graphicx,xspace,letltxmacro,xpatch}
\usepackage[english]{babel}
\usepackage{xfrac,xcolor}
\newtheorem{cdrthm}{Theorem}[section]
\newtheorem{theo}[cdrthm]{Theorem}
\newtheorem{prop}[cdrthm]{Proposition}
\newtheorem{lemm}[cdrthm]{Lemma}
\newtheorem{coro}[cdrthm]{Corollary}
\theoremstyle{definition}
\newtheorem{rema}[cdrthm]{Remark}
\newenvironment{proofc}{\begin{proof}[Proof (continued)]}{\end{proof}}


\newtheorem{defi}[cdrthm]{Definition}
\newtheorem*{sourceenoncestar}{\SourceEnonceName}
\NewDocumentEnvironment{enonce*}{O{plain} m o}{\def\SourceEnonceName{#2}\IfNoValueTF{#3}{\begin{sourceenoncestar}}{\begin{sourceenoncestar}[{#3}]}}{\end{sourceenoncestar}}



%
%\AtBeginDocument{
%%\newtheorem{myname}[cdrthm]{My beautiful theorem} }

\newcommand{\bfW}{\mathbf{W}}
\newcommand{\bbR}{\mathbb{R}}
\newcommand{\bbN}{\mathbb{N}}
\newcommand{\bbZ}{\mathbb{Z}}
\newcommand{\clF}{\mathcal{F}}
\newcommand{\clP}{\mathcal{P}}


\newcounter{steplcount}
\newenvironment{stepl}{\refstepcounter{steplcount}\begin{proof}[\textbf{Step~\thesteplcount}]}{\let\qed\relax\end{proof}}

%%%%%%%%%%%%%%%%%%%%%%%%%%%%%%%%%%%%%%%%%%%%%%%%%%%%%%%%%%%%%%%%%%
\graphicspath{{./figures/}}

\newcommand*{\mk}{\mkern -1mu}
\newcommand*{\Mk}{\mkern -2mu}
\newcommand*{\mK}{\mkern 1mu}
\newcommand*{\MK}{\mkern 2mu}

%\hypersetup{urlcolor=purple, linkcolor=blue, citecolor=red}

\newcommand*{\romanenumi}{\renewcommand*{\theenumi}{\roman{enumi}}}
\newcommand*{\Romanenumi}{\renewcommand*{\theenumi}{\Roman{enumi}}}
\newcommand*{\alphenumi}{\renewcommand*{\theenumi}{\alph{enumi}}}
\newcommand*{\Alphenumi}{\renewcommand*{\theenumi}{\Alph{enumi}}}
\let\oldtilde\tilde
\renewcommand*{\tilde}[1]{\mathchoice{\widetilde{#1}}{\widetilde{#1}}{\oldtilde{#1}}{\oldtilde{#1}}}
%%%%%%%%%%%%%%%%%%%%%%%%%%%%%%%%%%%%%%%%%%%%%%%%%%%%%%%%%%%%%%%%

\let\dpl\displaystyle
\providecommand{\skpt}{\par\smallskip\noindent}
\providecommand{\pointir}{.\ }
\providecommand{\eg}{e.g.}\providecommand{\ie}{i.e.}\providecommand{\etc}{etc.}\providecommand{\resp}{resp.}\providecommand{\cf}{cf.}
\providecommand{\mainmatter}{}\providecommand{\backmatter}{}
\let\Subsection\subsection\let\Subsubsection\subsubsection
\setlength{\emergencystretch}{2em}

\AtBeginDocument{\gdef\shorttitle{Geodesic stability without conjugate points}}
\title{Mañé structural stability characterizes Anosov geodesic flows in the original conjugate-free classes}
\author{Ludovic Rifford, Rafael Ruggiero}
\date{doi:10.5802/ahl.87; CC BY 4.0}
\begin{document}
\maketitle
\noindent\textit{Editorial note.} Complete original mathematical body and bibliography retained. Authors retain responsibility for mathematics.\par\medskip
\noindent\textit{Editorial note.} Original C2 structural stability from Mañé viewpoint, compact smooth no-conjugate-points setting and the precise dimension/quasi-convexity/diverging-rays/dense-generalized-rank-one alternatives are retained. No C1 closing lemma or arbitrary-dimensional unconditional characterization is added.\par\medskip
\noindent\textit{Allowed standard prerequisites (editorial).}\begin{itemize}
\item Standard spectral theory, functional analysis and original author-cited prior inputs. Original citations remain, and all source-local definitions, constructions and technical proofs are retained.
\end{itemize}\medskip
\section{Introduction}

The motivations for the main results in this article come from two sources. First of all, the challenging problem of the $C^{1}$ closing lemma for geodesic flows that remains an open, very difficult problem. Recently, Rifford~\cite{kn:Rifford} proved a $C^{0}$ closing lemma for geodesic flows applying ideas of geometric control theory. The high technical difficulties involved in the proof of this fact give an idea of the considerable complexity of the problem in the $C^{1}$ level. However, $C^{0}$ perturbations are considered too rough by specialists in perturbative theory of dynamical systems. So the question of how far we can go in proving the $C^{1}$ stability conjecture for geodesic flows without a $C^{1}$ closing lemma is an interesting, appealing problem in Riemannian geometry and dynamical systems.

Secondly, some important results about the topological dynamics of the geodesic flow of compact manifolds without conjugate points and hyperbolic global geometry are known for nonpositive curvature manifolds. Eberlein~\cite{kn:Eberlein} shows the topological transitivity of the geodesic flow of a visibility manifold, without restrictions on the sectional curvatures. The density of the set of periodic orbits, another important feature of topological dynamics, is known for visibility manifolds with nonpositive sectional curvatures (see for instance~\cite{kn:Ballmann}). Notice that by one of the main results of~\cite{kn:ARR}, the $C^{2}$ structural stability of the geodesic flow of a compact manifold from Mañé's viewpoint implies the hyperbolicity of the closure of the set of periodic orbits. So it seems natural to ask whether the density of periodic orbits, a statement with a flavor of topological dynamics, really needs some extra assumptions on the geometry of the manifold (like nonpositive curvature) to hold. There are some known results of course, Anosov geodesic flows have this property, as well as expansive geodesic flows in compact manifolds without conjugate points~\cite{kn:Ruggiero-expansive}. But visibility manifolds of nonpositive curvature are examples of the so-called rank one manifolds, which may have non Anosov geodesic flows because of the presence of flat strips.

The goal of the paper is to deal with the stability conjecture of geodesic flows of compact manifolds without conjugate points, a geometric condition that is much weaker than nonpositive curvature but still ensures many important properties for the topological dynamics of the geodesic flow.
%\bigskip
\goodbreak


\begin{theo}\label{main1}
Let $(M,g)$ be a compact $C^{\infty}$ manifold without conjugate points that is one of the following:
\begin{enumerate}
\item \label{theo1.1.1} A surface.
\item \label{theo1.1.2} A 3 dimensional manifold such that the universal covering is a quasi-convex space where geodesic rays diverge.
\end{enumerate}
Then, the geodesic flow is $C^{2}$ structurally stable from Mañé's viewpoint if and only if the geodesic flow is Anosov.
\end{theo}

%\bigskip

Item~\eqref{theo1.1.1} is probably known but we did not find any records in the literature about the subject. The second result for higher dimensional manifolds introduces a generalized version of the rank one notion for manifolds without conjugate points and no restrictions on the sectional curvatures (Section~1).


\begin{theo}\label{main2}
Let $(M,g)$ be a compact $C^{\infty}$ manifold without conjugate points such that the universal covering is a quasi-convex space where geodesic rays diverge. If the set of generalized rank one points is dense in $T_{1}M$ we have that the geodesic flow is $C^{2}$ structurally stable from Mañé's viewpoint if and only if the geodesic flow is Anosov.
\end{theo}

%\bigskip

The paper is organized as follows: Section~\ref{SECpreliminaries} is concerned with some preliminaries on manifolds without conjugate points. In Section~\ref{SECstrip}, we investigate the strip issue for manifolds without conjugate points. In Section~\ref{SECclosedorbits}, we study the stability and hyperbolicity properties of the set of closed orbits. Section~\ref{SECgromov} is concerned with the density of closed orbits on manifolds with Gromov hyperbolic fundamental group. Sections~\ref{sect.6} and~\ref{sect.7} are devoted to the proof of Theorem~\ref{main1} and the proof of Theorem~\ref{main2} is given in Section~\ref{SEChigher}.

\section{Preliminaries}\label{SECpreliminaries}

Let us give some notations and definitions that will be used through the article. A pair $(M,g)$ denotes a $C^{\infty}$ complete, connected, Riemannian manifold, $TM$ will denote its tangent space, $T_{1}M$ denotes its unit tangent bundle. $\Pi: TM \longrightarrow M$ denotes the canonical projection $\Pi(p,v) = p$, the coordinates $(p,v)$ for $TM$ will be called canonical coordinates. The universal covering of $M$ is $\tilde{M}$, the covering map is denoted by $\pi:\tilde{M} \longrightarrow M$, the pullback of the metric $g$ by $\pi$ is denoted by $\tilde{g}$. The geodesic $\gamma_{(p,v)}$ of $(M,g)$ or $(\tilde{M},\tilde{g})$ is the unique geodesic whose initial conditions are $\gamma_{(p,v)}(0)=p$, $\gamma_{(p,v)}'(0)= v$. All geodesics will be parametrized by arc length unless explicitly stated.

The fundamental group of $M$ will be denoted by $\pi_{1}(M)$.The group $\pi_{1}(M)$ acts by isometries in the universal covering $(\tilde{M}, \tilde{g})$. An element of the fundamental group when identified with an isometry of $(\tilde{M}, \tilde{g})$ will be called a \emph{covering isometry}.

\begin{defi}
A Riemannian manifold $(M,g)$ has no conjugate points if the exponential map is nonsingular at every point.
\end{defi}

Nonpositive curvature manifolds are well known examples of manifolds without conjugate points, but there are of course many examples of manifolds without conjugate points having sectional curvatures of variable sign. A particular family of such manifolds is the category of manifolds \emph{without focal points}: a geodesic of $(M,g)$ has no focal points if every nontrivial Jacobi field vanishing somewhere has increasing norm. The manifold $(M,g)$ has no focal points if every geodesic has no focal points.

The fundamental group of manifolds without conjugate points has some special algebraic properties similar to some properties of the fundamental group of manifolds with nonpositive curvature. For instance, $\pi_{1}(M)$ has no torsion if $(M,g)$ is compact and has no conjugate points. Given a covering isometry $h$, there exists a geodesic $\gamma_{\theta} \subset \tilde{M}$ such that
\[
h (\gamma_{\theta}(t)) = \gamma_{\theta}(t + P)
\]
for every $t \in \bbR$, where $P$ is a period of the periodic geodesic $\pi(\gamma_{\theta})$. Such a geodesic will be called an \emph{axis} of $(\tilde{M}, \tilde{g})$. The set of axes coincides with the set of lifts by the covering map of the closed geodesics of $(M,g)$. We shall denote by $A$ the set of axes of $(\tilde{M}, \tilde{g})$. For a more detailed description of the algebraic properties of the fundamental group of manifolds without conjugate points we refer to~\cite{kn:CS}.

\subsection{Divergence of geodesic rays in the universal covering, Busemann functions and invariant foliations}

We begin by recalling the following notion introduced by Eberlein~\cite{kn:Eberlein}.

\begin{defi}\label{divergence}
We say that geodesic rays \emph{diverge} in $(\tilde{M}, \tilde{g})$ if for given $p \in \tilde{M}$, $\epsilon >0$, $T>0$, there exists $R= R(p, \epsilon, T)>0$ such that two different geodesic rays $\gamma:\bbR \longrightarrow \tilde{M}$, $\beta : \bbR \longrightarrow \tilde{M}$ with $\gamma(0)= \beta(0) =p$, subtending an angle at $p$ greater than $\epsilon$, then $d(\gamma(t),\beta(t)) \geq T$ for every $t \geq R$. We say that geodesic rays \emph{diverge uniformly} in $(\tilde{M}, \tilde{g})$ if the number $R(p, \epsilon, T)$ is independent of $p \in \tilde{M}$.
\end{defi}

The divergence of geodesic rays is common to all known categories of manifolds without conjugate points (nonpositive curvature, no focal points, bounded asymptote~\cite{kn:Eschenburg}, compact surfaces without conjugate points), but it is a conjecture whether it is satisfied for every compact manifold without conjugate points. There are some partial results pointing towards a positive answer (see for instance~\cite{kn:Ruggiero-div}).

One of the main basic features of manifolds without conjugate points is the existence of many asymptotic objects which describe the global geometry of $(\tilde{M}, \tilde{g})$. We start by recalling the notion of Busemann function.

\begin{defi}\label{busemann}
Given $\theta = (p,v) \in T_{1}\tilde{M}$, the Busemann function $b^{\theta} : \tilde{M} \longrightarrow \bbR$ is given by
\[
b^{\theta}(x) = \lim_{t\rightarrow +\infty} (d(x,\gamma_{\theta}(t)) -t).
\]
\end{defi}

Busemann functions of compact manifolds without conjugate points are $C^{1+k}$, namely, $C^{1}$ functions with $k$-Lipschitz first derivatives where the constant $k$ depends on the minimum value of the sectional curvatures (see for instance~\cite[Section~6]{kn:Pesin}). The level set
\[
(b^{\theta})^{-1}(-t)= H_{\theta}(t)
\]
is called \emph{horosphere}. For each $\theta \in T_{1}\tilde{M}$, the collection of sets $H_{\theta}(t)$, $t \in \bbR$, defines a $C^{1}$ foliation of $\tilde{M}$ by equidistant leaves. The gradient $\nabla b^{\theta}$ is a Lipschitz unit vector field and its flow $\psi^{\theta}_{t}: \tilde{M} \longrightarrow \tilde{M}$ that is always tangent to geodesics of $(\tilde{M}, \tilde{g})$ preserves the foliation of the horospheres, namely
\[
\psi^{\theta}_{s}(H_{\theta}(t)) = H_{\theta}(t-s)
\]
for every $t, s \in \bbR$. The integral orbits of the Busemann flow are usually called \emph{Busemann asymptotes} of $\gamma_{\theta}$. When the curvature is nonpositive, Busemann functions and horospheres are $C^{2}$ smooth~\cite{kn:Eschenburg}.

\begin{defi}\label{centerstable}
For $\theta = (p,v) \in T_{1}\tilde{M}$ let
\begin{align*}
\tilde{W}^{s}(\theta) &= \left\{\left(x, -\nabla_{x}b^{\theta}\right), \, x \in H_{\theta}(0) \right\}
\\
\tilde{W}^{u}(\theta) &= \left\{\left(x, \nabla_{x}b^{(p,\,-v)}\right), \, x \in H_{(p,-v)}(0)\right\}.
\end{align*}

We shall denote by $\widetilde{\clF}^{s}$ the collection of the sets $\tilde{W}^{s}(\theta)$, and by $\widetilde{\clF}^{u}$ the collection of the sets $\tilde{}W^{u}(\theta)$, $\theta \in T_{1}\tilde{M}$. If $\clP : T_{1}\tilde{M} \longrightarrow T_{1}M$ is projection $\clP(p,v) = (\pi(p), d_{p}\pi(v))$, let
\[
W^{s}(\theta) = \clP(\tilde{W}^{s}(\theta))
\]
\[
W^{u}(\theta) = \clP(\tilde{W}^{u}(\theta)).
\]
Let us denote by $\clF^{s}$ the collection of the sets $W^{s}(\theta)$, $\theta \in T_{1}M$, and by $\clF^{u}$ the collection of the sets $W^{u}(\theta)$, $\theta \in T_{1}M$.
\end{defi}

The sets $W^{s}(\theta)$, $W^{u}(\theta)$, are Lipschitz continuous, $(m-1)$ dimensional submanifolds of $T_{1}M$, where $m= dim(M)$, and of course they coincide with the stable and the unstable sets of $\theta$ when the geodesic flow is Anosov. When $(M,g)$ is a compact surface without conjugate points, the collections $\clF^{s}$, $\clF^{u}$, form continuous foliations. This is a well known consequence of the divergence of geodesic rays in the universal covering proved by Green~\cite{kn:Green} and the quasi-convexity of $(\tilde{M}, \tilde{g})$ shown by Morse~\cite{kn:Morse}. The same holds for any compact manifold with nonpositive curvature, no focal points, and for a rather more general category of manifolds satisfying a condition called \emph{bounded asymptote} (see~\cite{kn:Eschenburg}).

The most general known result concerning the regularity of the above family of invariant sets is proved in~\cite{kn:Ruggiero-Asterisque}.

\begin{theo}\label{divergence-continuity}
Let $(M,g)$ be a compact manifold without conjugate points. Then geodesic rays diverge uniformly in $\tilde{M}$ if and only if $\clF^{s}$, $\clF^{u}$, $\widetilde{\clF}^{s}$, $\widetilde{\clF}^{u}$ are continuous foliations by Lipschitz leaves invariant by the geodesic flow (in $(M,g)$, $(\tilde{M}, \tilde{g})$ respectively).
\end{theo}

\subsection{Quasi-convexity}

\begin{defi}\label{quasi-convexity}
The universal covering $(\tilde{M}, \tilde{g})$ of a complete Riemannian manifold $(M,g)$ is called a $(K,C)$-quasi-convex space, or simply a quasi-convex space, if there exist constants $K>0$, $C>0$, such that given two pairs of points $x_{1}, x_{2}$, $y_{1},y_{2}$ in $\tilde{M}$ and two minimizing geodesics $\gamma:[0,1] \longrightarrow \tilde{M}$, $\beta:[0,1] \longrightarrow \tilde{M}$ such that $\gamma(0)= x_{1}$, $\gamma(1)= y_{1}$, $\beta(0) = x_{2}$, $\beta(1)= y_{2}$, we have
\[
d_{H}(\gamma, \beta) \leq K\max\{ d(x_{1},x_{2}), d(y_{1},y_{2})\} + C
\]
where $d_{H}$ is the Hausdorff distance.
\end{defi}

The universal covering of manifold of nonpositive sectional curvature is $(1,0)$-quasi-convex. Most of the known categories of manifolds without conjugate points (no focal points, bounded asymptote) have quasi-convex universal coverings. Moreover, by the work of Gromov~\cite{kn:Gromov}, the universal covering of every compact manifold whose fundamental group is hyperbolic is quasi-convex. Although geodesics in $(\tilde{M},\tilde{g})$ behave like hyperbolic geodesics when the dimension of $M$ is $2$, an example by Ballmann--Brin--Burns~\cite{kn:BBB} shows that Jacobi fields may behave wildly compared with the quasi-convex behavior of geodesics.

The link between quasi-convexity and the regularity of the families of sets $\clF^{s}$, $\clF^{u}$ is the following result proved in~\cite{kn:Ruggiero-Ensaios}.

\begin{theo}\label{quasi-conv-continuity}
Let $(M,g)$ be a compact manifold without conjugate points such that $(\tilde{M}, \tilde{g})$ is quasi-convex and geodesic rays diverge. Then geodesic rays diverge uniformly.
\end{theo}

Theorems~\ref{divergence-continuity} and~\ref{quasi-conv-continuity} then imply that the families $\clF^{s}$, $\clF^{u}$ form continuous invariant foliations whenever $(\tilde{M}, \tilde{g})$ is quasi-convex and geodesic rays diverge.


\subsection{Busemann asymptotes versus asymptotes}

As we mentioned before, every orbit of the Busemann flow of $b^{\theta}$ is called a Busemann asymptote of $\gamma_{\theta}$ for $\theta \in T_{1}\tilde{M}$. However, the usual definition of asymptoticity is the following:

\begin{defi}
A geodesic $\beta \subset \tilde{M}$ is forward asymptotic to $\gamma \subset \tilde{M}$ if there exists $L>0$ such that
\[
d_{H}\left(\gamma[0,+\infty), \beta[0,+\infty)\right) \leq L
\]
where $\gamma[0,+\infty) = \{ \gamma(t), t \geq 0\}$. A geodesic $\sigma \subset \tilde{M}$ is backward asymptotic to $\gamma$ if there exists $L>0$ such that
\[
d_{H}\left(\gamma(-\infty,0], \beta(-\infty,0]\right) \leq L.
\]
Two geodesics $\gamma$, $\beta$ in $\tilde{M}$ are bi-asymptotic if they are both forward and backward asymptotic.
\end{defi}

The quasi-convexity of $(\tilde{M}, \tilde{g})$ implies that the global behavior of Busemann asymptotes is a coarse version of the behavior of asymptotes in nonpositive curvature (see~\cite{kn:Ruggiero-Ensaios} for instance):

\begin{lemm}\label{asymptotic-quasi-convex}
Let $(M,g)$ be a compact manifold without conjugate points such that $(\tilde{M}, \tilde{g})$ is $K,C$ quasi-convex. Then for every $\theta \in T_{1}\tilde{M}$, $\theta= (p,v)$, and every $(q, w) \in H_{\theta}(0)$, we have
\[
d\left(\gamma_{\theta}(t), \gamma_{(q,\,w)}(t)\right) \leq Kd(p,q) +C
\]
for every $ t \geq 0$.
\end{lemm}

However, a Busemann asymptote of $\gamma_{\theta}$ might not be asymptotic to $\gamma_{\theta}$. What is true is

\begin{lemm}\label{asymptotic-1}
Let $(M,g)$ be compact without conjugate points such that geodesic rays diverge uniformly in $\tilde{M}$. Then
\begin{enumerate}
\item \label{lemm2.10.1}Any geodesic $\gamma_{\eta}$ forward asymptotic to $\gamma_{\theta}$ is a Busemann asymptote of $\gamma_{\theta}$. Moreover, there exists $c>0$ such that $b^{\eta} (x) = b^{\theta}(x) +c$ for every $ x \in \tilde{M}$.
\item  \label{lemm2.10.2}In particular, any geodesic $\gamma_{\eta}$ that is bi-asymptotic to $\gamma_{\theta}$ is a Busemann asymptote of $\gamma_{\theta}$ and $\gamma_{-\theta}$ where $\theta = (p,v)$, $-\theta= (p,-v)$. In this case, if $\gamma_{\eta}(0) \in H_{\theta}(0)$ then $\gamma_{\eta}(0) \in H_{\theta}(0)\cap H_{-\theta}(0)$.
\end{enumerate}
\end{lemm}

\begin{proof}
We follow the ideas in~\cite{kn:Ruggiero-Asterisque}, most of them are well known to specialists.

Item~\eqref{lemm2.10.1} follows from the fact that the horosphere $H_{\theta}(0)$ is the limit of spheres $S_{t}(\gamma_{\theta}(t))$ of radius $t$ centered at $\gamma_{\theta}(t)$. Indeed, let $ x \in \tilde{M}$ and consider the geodesics $[x, \gamma_{\theta}(t)]$ joining $x$ and $\gamma_{\theta}(t)$, and $[x,\gamma_{\eta}(t)]$ joining $x$ and $\gamma_{\eta}(t)$. The angle subtended by these geodesics at $x$ tends to zero as $t \rightarrow +\infty$ because of the uniform divergence of geodesic rays. This yields that $\gamma_{\eta}$ and $\gamma_{\theta}$ have the same Busemann asymptotes. Since the Busemann asymptotes of $\gamma_{\theta}$ define a flow by geodesics that is always orthogonal to the horospheres $H_{\theta}(s)$, we have that the horospheres $H_{\theta}(s)$ and $H_{\eta}(r)$ give rise to two foliations which are perpendicular to the same flow by geodesics. Hence the functions $b^{\eta}$ and $\beta^{\theta}$ have the same gradients and since they are $C^{1}$ they differ by a constant. This proves item~\eqref{lemm2.10.1}.

To show item~\eqref{lemm2.10.2} we have to prove that if $\gamma_{\eta}(0) \in H_{\theta}(0)$ and $\gamma_{\eta}$ is bi-asymptotic to $\gamma_{\theta}$, then $\gamma_{\eta}(0)$ belongs to $H_{-\theta}(0)$ as well. Item (1) implies that, since $\gamma_{-\eta}$ is asymptotic to $\gamma_{-\theta}$, then both have the same Busemann flows and their Busemann functions differ by a constant. This yields that there exists $a>0$ such that
\[
\gamma_{\eta}(0) = \gamma_{-\eta}(0) \in H_{-\theta}(a),
\]
and hence $H_{-\eta}(0) = H_{-\theta}(a)$. We must show that $a =0$.

Let $\theta = (p,v)$, $\eta = (q, w)$, we know that $H_{\theta}(0)$ and $H_{-\theta}(0)$ are tangent at $p$, and that $\tilde{M}- H_{\theta}(0)$ consists of two disjoint open regions
\begin{align*}
O_{\theta}^{+} &= \left\{x\in \tilde{M}, \,b^{\theta}(x) >0\right\}
\\
O_{\theta}^{-} &= \left\{x\in \tilde{M}, \,b^{\theta}(x) <0\right\}.
\end{align*}

Moreover we have
\[
O_{-\theta}^{-} \subset O_{\theta}^{+},
\]
\[
H_{-\theta}(0) \subset O_{\theta}^{+} \cup \{p \},
\]
\[
H_{-\eta}(0) \subset O_{\theta}^{+} \cup \{q\}
\]
since $H_{\eta}(0) = H_{\theta}(0)$.

Assume that $a>0$. Then $b^{-\theta}(H_{-\theta}(a)) =-a$ and therefore $H_{-\theta}(a) \subset O_{-\theta}^{-}$, that is strictly contained in $O_{\theta}^{+}$. This contradicts the fact that $q \in H_{-\theta}(a) \cap \{ (b^{\theta})^{-1}(0) \}$.

If on the other hand, $a<0$, then the horosphere $H_{-\theta}(a)$ intersects the open set $O_{\theta}^{-}$ which is again a contradiction since $H_{-\theta}(a) = H_{-\eta}(0) \subset O_{\theta}^{+} \cup \{q\}$. So $a=0$ thus finishing the proof of the Lemma.
\end{proof}

\pagebreak
Combining the above two Lemmas~\ref{asymptotic-quasi-convex} and~\ref{asymptotic-1} we get

\begin{coro}\label{asympt-unique}
Let $(M,g)$ be a compact manifold without conjugate points whose universal covering is quasi-convex where geodesic rays diverge. Then two geodesics of $(\tilde{M}, \tilde{g})$ are asymptotic if and only if they are Busemann asymptotic (up to reparametrization).
\end{coro}

A natural question arises from the previous Lemma: Is there any connected set of bi-asymptotic geodesics containing two given bi-asymptotic geodesics? This would be in many respects an analogous to the flat strip theorem. What we know is the following result proved by Croke--Schroeder~\cite{kn:CS}.

\begin{theo}\label{CS}
Let $(M,g)$ be a compact analytic manifold without conjugate points. Then the set of closed geodesics in a given nontrivial homotopy class is a connected, rectifiable set (each pair of points can be joined by a rectifiable curve in the set) of closed geodesics in the same homotopy class and constant period.
\end{theo}

\subsection{Generalized rank one manifolds}

Let us recall that a manifold with nonpositive curvature is called a \emph{rank one manifold} if there exists a geodesic where the dimension of parallel Jacobi fields along this geodesic is one. Every compact surface with nonpositive curvature and genus greater than one is a rank one manifold of course.

Theorem~\ref{divergence-continuity} leads naturally to the following extension of the notion of rank one manifold.

\begin{defi}\label{DEFgenrankone}
A generalized rank one manifold is a compact manifold without conjugate points such that
\begin{enumerate}
\item  \label{defi2.13.1}Geodesic rays diverge uniformly in $(\tilde{M}, \tilde{g})$.
\item \label{defi2.13.2}There exists $\theta \in T_{1}M$ and an open neighborhood $B(\theta)$ of $\theta$ in $T_1M$ such that for each $\eta \in B(\theta)$, the connected component $W^{s}_{loc}(\eta)$ of $W^{s}(\eta) \cap B(\eta)$ containing $\eta$ and the connected component $W^{u}_{loc}(\eta)$ of $W^{u}(\eta) \cap B(\eta)$ containing $\eta$ satisfy
\[
W^{s}_{loc}(\eta) \cap W^{u}_{loc}(\eta) = \eta.
\]
The point $\theta$ will be called a generalized rank one point for the geodesic flow.
\end{enumerate}
\end{defi}

Notice that by definition, if the set of generalized rank one points is non empty then it is open. It is clear that rank one manifolds of nonpositive curvature are generalized rank one manifolds since geodesic rays diverge uniformly in nonpositive curvature and the tangent space of $ W^{s}_{loc}(\eta) \cap W^{u}_{loc}(\eta)$ is generated by parallel Jacobi fields which are linearly independent from the geodesic vector field $\gamma_{\eta}'(0)$. So along the orbit of a rank one point in a space of nonpositive curvature stable and unstable sets intersect transversally, and since these sets form continuous foliations by $C^{1}$ leaves the transversality between invariant submanifolds is an open property.

However, the set of generalized rank one points might include strictly the set of rank one points of manifolds with nonpositive curvature. This is the case of surfaces of nonpositive curvature where the curvature vanishes just at a finite set of closed geodesics. Every point of $T_{1}M$ is a generalized rank one point while the set of rank one points is the complement of this finite set of flat geodesics. The expansivity of the geodesic flows of a compact manifold without conjugate points implies that every point in $T_{1}M$ is a generalized rank one point (see~\cite{kn:Ruggiero-expansive}). We might expect that the presence of generalized rank one points would imply some kind of local expansivity, this will be the subject of the last section.

The study of the set of intersections $\tilde{W}^{s}(\eta) \cap \tilde{W}^{u}(\eta)$ is one of the most intriguing problems in the theory of manifolds without conjugate points. In the case of compact surfaces such a set is a connected compact curve with boundary (that might be a single point of course). Manifolds without focal points satisfy the so-called \emph{flat strip Theorem}: every two bi-asymptotic geodesics in $\tilde{M}$ bound a flat, isometrically embedded strip $[0,a]\times \bbR$ in $\tilde{M}$.

The convexity properties of spaces of nonpositive curvature yield that the set of bi-asymptotic geodesics to a given one is a convex flat set (see for instance~\cite{kn:BGS} and references about the subject)

Without restrictions on the sectional curvatures the intersections between invariant submanifolds might be non-flat strips as shown by Burns (see~\cite{kn:Burns}), but still enjoy good topological properties. In higher dimensions this problem is much more difficult, this will be the subject of the next section.


\section{The strip issue for manifolds without conjugate points}\label{SECstrip}

If we drop any assumption on the sectional curvatures or Jacobi fields, or even the analytic hypothesis considered by Croke and Schroeder, we can show the following result about the topology of the set of bi-asymptotic geodesics that is new in the theory and interesting in itself.

\begin{lemm}\label{strip}
Let $(M,g)$ be a compact $C^{\infty}$ manifold without conjugate points such that $(\tilde{M}, \tilde{g})$ is $(K,C)$-quasi-convex where geodesic rays diverge. Then given $\theta = (p,v) \in T_{1}\tilde{M}$, and a geodesic $\beta=\gamma_{\eta}$ (with $\eta \in T_{1}\tilde{M}$) bi-asymptotic to $\gamma = \gamma_{\theta}$, there exists a connected set $\Sigma(\gamma, \beta) \subset H_{\theta}(0) \cap H_{-\theta}(0)$ containing $p$ and $\beta \cap H_{\theta}(0)$, such that for every $x \in \Sigma(\gamma, \beta)$, the geodesic with initial conditions $(x, -\nabla_{x} b^{\theta})$ is bi-asymptotic to both of them. In particular, the set
\[
S(\gamma, \beta) = \bigcup_{x\,\in\,\Sigma(\gamma,\,\beta),\,t\, \in\,\bbR} \gamma_{\left(x,\,-\nabla_{x}\,b^{\theta}\right)}(t)
\]
is homeomorphic to $\Sigma(\gamma, \beta) \times \bbR$.
\end{lemm}

Before proving Lemma~\ref{strip}, let us demonstrate the following elementary lemma.

\begin{lemm}\label{lemLR}
Let $(M,g)$ be a compact $C^{\infty}$ manifold without conjugate points and $\theta = (p,v) \in T_{1}\tilde{M}$ be fixed. Then for every $x\in \tilde{M}$, any $u\in \bbR$ such that
\[
d(x,\gamma) := \inf_{t\,\in\,\bbR} \left\{d \left(x,\gamma(t) \right) \right\} = d(x,\gamma(u))
\]
satisfies
\[
\left| u + b^{\theta}(x)\right| \leq d(x,\gamma).
\]
\end{lemm}

\begin{proof}
Let $x\in \tilde{M}$ and $u\in \bbR $ be such that $d(x,\gamma) =d(x,\gamma(u))$, for every $t>0$ the triangle inequality yields
\begin{align*}
d\left(x,\gamma(t) \right) &\leq d\left(x,\gamma(u)\right) + d\left(\gamma(u),\gamma(t)\right) = d(x,\gamma) + |t-u|\\
\intertext{and}
|t-u| = d\left(\gamma(u),\gamma(t) \right) &\leq d\left(\gamma(u),x\right) + d\left(x,\gamma(t)\right) = d(x,\gamma) + d\left(x,\gamma(t)\right).
\end{align*}
We conclude easily by letting $t$ tend to $+\infty$ and using the definition of $b^{\theta}(x)$.
\end{proof}

\begin{proof}[Proof of Lemma~\ref{strip}]
We construct the set $\Sigma(\gamma, \beta)$ by hand. For every $t\in \bbR$, let $c_{t}: [0,1] \longrightarrow \tilde{M}$ be the geodesic with $c_{t}(0)= \gamma_{\theta}(t)$, $c_{t}(1) = \beta(t)$. For every positive integer $n$ and every $s\in [0,1]$, we consider the geodesic segment $\alpha_n^s : [0,1] \rightarrow \tilde{M}$ joining $c_{-n}(s)$ to $c_{n}(s)$. Since geodesic rays diverge in $\tilde{M}$, by Lemma~\ref{asymptotic-1}, there is $c\in \bbR$ such that $b^{\eta}=b^{\theta}+c$. Thus, for every integer $n$ if $v, w \in \bbR$ satisfy
\[
d \left(\beta(n),\gamma\right) = d \left(\beta(n),\gamma(v) \right) \quad \mbox{and} \quad d \left(\beta(-n),\gamma\right) = d \left(\beta(-n),\gamma(w) \right)
\]
then by Lemma~\ref{lemLR} there holds
\begin{align*}
\left| v-n-c\right| = \left|v+b^{\eta}(\beta(n)) -c \right| &= \left|v+b^{\theta}(\beta(n))\right| \leq d \left(\beta(n),\gamma\right) \leq d_{H} \left(\gamma,\beta\right)
\\
\intertext{and}
\left| v+n-c\right| = \left|v+b^{\eta}(\beta(-n)) -c \right| &= \left|v+b^{\theta}(\beta(-n))\right| \leq d \left(\beta(-n),\gamma\right) \leq d_{H} \left(\gamma,\beta\right),
\end{align*}
which shows that
\begin{align*}
d\left(\beta(n),\gamma(n)\right) & \leq d\left(\beta(n),\gamma(v)\right) + d\left(\gamma(v),\gamma(n)\right)\\
& \leq d\left(\beta(n),\gamma\right) + |v-n| \\
& \leq 2 d_{H} \left(\gamma,\beta\right) + |c|,
\end{align*}
and in the same way that $d\left(\beta(-n),\gamma(-n)\right) \leq 2 d_{H} \left(\gamma,\beta\right) + |c|$. Let $D:=d_H(\gamma,\beta)$ and fix an integer $n>2D+|c|$ and $s\in [0,1]$, we have
\begin{align*}
d\left(\alpha_n^s(1),\gamma(n)\right) &\leq d\left(\beta(n),\gamma(n)\right) \leq 2 d_{H} \left(\gamma,\beta\right) + |c|
\\
\intertext{and}
d\left(\alpha_n^s(0),\gamma(-n)\right) &\leq d\left(\beta(-n),\gamma(-n)\right) \leq 2 d_{H} \left(\gamma,\beta\right) + |c|.
\end{align*}
So that for every $t>0$
\begin{align*}
d\left(\alpha_n^s(0),\gamma(t)\right) & \geq d\left(\gamma(t),\gamma(-n)\right)-d\left(\alpha_n^s(0),\gamma(-n)\right) \\
& \geq t+n -2 D -|c| >0
\\
\intertext{and}
d\left(\alpha_n^s(1),\gamma(t)\right) & \leq d\left(\gamma(t),\gamma(n)\right)+d\left(\alpha_n^s(1),\gamma(n)\right) \\
& \leq t-n +2 D +|c| <0.
\end{align*}
Taking the limit as $t$ tends to $+\infty$, we infer that $b^{\theta}(\alpha_n^s(0))>0$ and $b^{\theta}(\alpha_n^s(1))<0$. As a consequence, there is $r\in (0,1)$ such that $\alpha_n^s(r)$ belongs to $H_{\theta}(0)$. By Lemma~\ref{lemLR}, there is $u\in \bbR$ such that
\[
d \left(\alpha_n^s(r),\gamma\right) = d \left(\alpha_n^s(r),\gamma(u) \right) \quad \text{and} \quad |u| \leq d \left(\alpha_n^s(r),\gamma\right),
\]
which by quasi-convexity together with the above inequalities gives
\begin{align*}
|u| \leq d_{H}\left(\alpha_n^s, \gamma\right) & \leq K\max \left\{ d\left(c_{-n}(s),\gamma(-n)\right), d\left(c_n(s),\gamma(n)\right)\right\} + C\\
& \leq K\max \left\{ d\left(\beta(-n),\gamma(-n)\right), d\left(\beta(n),\gamma(n)\right)\right\} + C \\
& \leq 2KD + K|c| +C,
\end{align*}
and in turn
\begin{align*}
d \left(\alpha_n^s(r),p \right) & \leq d \left(\alpha_n^s(r),\gamma(u)\right) + d\left(\gamma(u),\gamma(0)\right) \\
& = d \left(\alpha_n^s(r),\gamma\right) + |u| \\
& \leq 4KD + 2K|c| +2C =:\tau.
\end{align*}
By the divergence of rays in $\tilde{M}$, the geodesics $[c_{-n}(s), c_{n}(s)]$ tend to be orthogonal to $H_{\theta}(0)$ at their points of intersection. So for $n$ large, there is a unique $r_{n}^s \in (0,1)$ such that $\alpha_n^s(r_n^s) \in H_{\theta}(0)$ and the mapping
\[
\Gamma_n \, : \, s \in [0,1] \longmapsto \alpha_n^s(r_n^s) \in H_{\theta}(0) \cap B_{\tau}(p)
\]
is continuous (here $B_{\tau}(p)$ stands for the closed ball centered at $p$ with radius $\tau$). Let $\Sigma(\gamma, \beta)$ be the set of $q\in \tilde{M}$ for which there exists a sequence $\{n_k\}_k$ of positive integers tending to infinity such that
\[
q = \lim_{k\rightarrow \infty} \Gamma_{n_k}\left(s_{n_k}\right).
\]
By construction, $\Sigma(\gamma, \beta)$ is a closed subset of $H_{\theta}(0)$ contained in $B_{\tau}(p)$, which contains $p=\gamma(0)$ and $q :=\beta \cap H_{\theta}(0)$. We claim that $\Sigma(\gamma, \beta)$ is connected. As a matter of fact, if there are two disjoint open subsets $A_1, A_2$ of $H_{\theta}(0)$ such that $\Sigma(\gamma, \beta)\subset A_1 \cup A_2$ with $p\in A_1$, then all points of $\Sigma(\gamma,\beta)$ must belong to $A_1$ because otherwise there is a sequence of continuous curves in $H_{\theta}(0)\cap B_{\tau}(p)$, given by the restrictions of some $\Gamma_n$, which connects $p$ to $A_2$ that gives rise, by compactness, to an accumulation point in $\Sigma(\gamma,\beta)$ outside $A_1\cup A_2$, a contradiction.

To finish the proof of the Lemma~\ref{strip}, it remains to show that $\Sigma(\gamma,
\beta)$ is a subset of $H_{\theta}(0)\cap H_{-\theta}(0)$.

Observe first that by quasi-convexity, for every integer $n$ and every $s\in [0,1]$,
\[
d_{H}\left(\gamma,\alpha_n^s\right) \leq 2KD + K|c| +C,
\]
so any convergent subsequence of the points $\alpha_{n_{k}}(s_{n_{k}}) \rightarrow q \in H_{\theta}(0)$ gives rise to a geodesic $\sigma^{q}(t)$ satisfying
\[
d_{H}(\gamma_{\theta}, \sigma^{q}) \leq K d_{H}(\gamma_{\theta}, \beta) +C,
\]
meaning that $\sigma^{q}$ is bi-asymptotic to $\gamma_{\theta}$. Moreover, the assumptions on the lemma allow to apply Theorem~\ref{quasi-conv-continuity}: geodesic rays diverge uniformly in $\tilde{M}$. Therefore, Lemma~\ref{asymptotic-1} proceeds and we get $\Sigma(\gamma, \beta) \subset H_{\theta}(0)\cap H_{-\theta}(0)$.
\end{proof}

\begin{coro}\label{connected-strip}
Let $(M,g)$ be a compact manifold without conjugate points such that $(\tilde{M}, \tilde{g})$ is quasi-convex where geodesic rays diverge. Then given a geodesic $\gamma_{\theta} \subset \tilde{M}$ the set $S(\gamma_{\theta})$ of geodesics which are bi-asymptotic to $\gamma_{\theta}$ is homeomorphic to a product $\Sigma(\gamma_{\theta}) \times \bbR$ where $\Sigma(\gamma_{\theta}) \subset H_{\theta}(0)\cap H_{-\theta}(0)$ is a connected set.
\end{coro}

\begin{proof}
This is straightforward from Lemma~\ref{strip}.
\end{proof}

Some remarks about Corollary~\ref{connected-strip}. The set $S(\gamma)$ is a natural candidate to be ``the strip'' of $\gamma$. However, although it is homeomorphic to a product of the line and a compact connected set like in nonpositive curvature, its geometry might be quite wild.

\section{Stability and hyperbolicity of the set of closed orbits}\label{SECclosedorbits}

In this section we remind some of the main steps of the proof of the stability conjecture for diffeomorphisms that can be extended to geodesic flows, notably after a recent version of the Franks' Lemma for the so-called Mañé perturbations of a Riemannian metric.

We start by recalling some basic definitions concerning hyperbolic dynamics.

\begin{defi}\label{hyperbolic1}
Let $\psi_{t}: N\longrightarrow N$ be a smooth flow without singularities acting on a complete $C^{\infty}$ Riemannian manifold. An invariant set $Y \subset N$ is called hyperbolic if there exists $C>0$, $r>0$, and for every $ p \in Y$ there exist subspaces $E^{s}(p)$, $E^{u}(p)$ such that
\begin{enumerate}
\item \label{defi4.1.1}$E^{s}(p) \oplus E^{u}(p) \oplus X(p) = T_{p}N$ where $X(p)$ is the subspace tangent to the flow.
\item \label{defi4.1.2}$\parallel d_{p}\psi_{t}(v) \parallel \leq Ce^{-rt} \parallel v \parallel $ for every $t \geq 0$ and $v \in E^{s}(p)$.
\item \label{defi4.1.3}$\parallel d_{p}\psi_{t}(v) \parallel \leq Ce^{rt} \parallel v \parallel $ for every $t \leq 0$ and $v \in E^{u}(p)$.
\end{enumerate}

The subspace $E^{s}(p)$ is called stable subspace, the subspace $E^{u}(p)$ is called the unstable subspace. When $Y = N$ the flow $\psi_{t}$ is called Anosov.

Replacing $\psi_{t}$ by $f^{t}$, where $f:N \longrightarrow N$ is a diffeomorphism, $t \in \bbN$, and erasing $X(p)$ from item~\eqref{defi4.1.1} in the above definition, we get what is called a hyperbolic set for the diffeomorphism $f$. If $Y=N$ the diffeomorphism $f$ is called an Anosov diffeomorphism.
\end{defi}

The theory of hyperbolic sets of flows and diffeomorphisms is very rich, one of the main features of the dynamics is the existence of invariant submanifolds $\bfW^{s}(p)$, $\bfW^{u}(p)$ for every $p$ in the hyperbolic set where asymptotic properties of orbits are counterparts of asymptotic properties of the differential of the system acting on stable and unstable subspaces (see for instance~\cite{kn:Anosov,kn:HPS}). The submanifold $\bfW^{s}(p)$ is always tangent to the bundle $E^{s}$, the submanifold $\bfW^{u}(p)$ is always tangent to the bundle $E^{u}$.

By Lemma~\ref{asymptotic-1}, the invariant submanifolds $\bfW^{s}(\theta)$, $\bfW^{u}(\theta)$ agree with the sets $W^{s}(\theta)$, $W^{u}(\theta)$ locally when the geodesic flow of a manifold without conjugate points is Anosov. This fact together with Corollary~\ref{connected-strip} imply

\begin{lemm}\label{hyperbolic-no-strip}
Let $(M,g)$ be a compact manifold without conjugate points such that $(\tilde{M}, \tilde{g})$ is quasi-convex and geodesic rays diverge. If the orbit of $\theta \in T_{1}M$ is a hyperbolic periodic orbit then for every lift $\gamma_{\eta} \subset \tilde{M}$ of the geodesic $\gamma_{\theta}$ we have
\[
S(\gamma_{\eta}) = \gamma_{\eta}.
\]
\end{lemm}

\begin{proof}
The proof is straightforward from Corollary~\ref{connected-strip}: The set $S(\gamma_{\eta})$ contains the intersection $\Sigma(\gamma_{\eta}) = H_{\eta}(0) \cap H_{-\eta}(0)$, which, according to the definition of $\tilde{W}^{s}(\eta)$, $\tilde{W}^{u}(\eta)$, is homeomorphic to $\tilde{W}^{s}(\eta) \cap \tilde{W}^{u}(\eta)$ (with respect to the relative topology of the sets). By the hyperbolicity of the periodic orbit and Lemma~\ref{asymptotic-1}, we have that
\[
\tilde{W}^{s}(\eta) \cap \tilde{W}^{u}(\eta) = \bfW^{s}(\eta) \cap \bfW^{u}(\eta) = \eta
\]
which yields the statement.
\end{proof}


Actually, if the geodesic flow of a compact Riemannian manifold is Anosov then the manifold has no conjugate points by a celebrated theorem due to Klingenberg~\cite{kn:Klingenberg} (see also a nice generalization by R. Mañé~\cite{kn:Mane2}).

Systems with hyperbolic invariant sets are closely related to the theory of stable systems.

\begin{defi}\label{Ck-stability}
A smooth flow $\psi_{t} : N\longrightarrow N$ acting on a smooth manifold is $C^{k}$ structurally stable if there exists $\epsilon >0$ such that every flow $\rho_{t}$ in the $\epsilon$-neighborhood of $\psi_{t}$ in the $C^{k}$ topology is conjugate to $\psi_{t}$. Namely, there exists a homeomorphism $h_{\rho}: N\longrightarrow N$ such that
\[
h(\psi_{t}(p)) = \rho_{s_{p}(t)}(h(p))
\]
for every $t \in \bbR$, where $s_{p}(t)$ is a continuous injective function with $s_{p}(0) =0$.
\end{defi}

A series of results in the 60's, 70's and 80's characterize $C^{1}$ structurally stable systems (mainly~\cite{kn:Mane,kn:Mane88,kn:Robin-Robinson,kn:Robin-Robinson2}).

\begin{theo}\label{stab-diffeo}
A diffeomorphism acting on a compact manifold is $C^{1}$ structurally stable if and only if it is Axiom A and the intersections between stable and unstable manifolds of periodic orbits is always transversal.
\end{theo}

Recall that a diffeomorphism $f$ acting on a smooth manifold is called \emph{Axiom A} if the set $\Omega(f)$ of nonwandering points of $f$ is a hyperbolic set and the set of periodic orbits is dense in $\Omega(f)$.

Theorem~\ref{stab-diffeo} characterizes as well $C^{1}$ structurally stable flows without singularities on compact manifolds. Newhouse~\cite{kn:Newhouse} shows that a symplectic diffeomorphism acting on a compact manifold is $C^{1}$ structurally stable if and only if it is Anosov. Newhouse's proof extends to energy levels without singularities of Hamiltonian flows defined on the cotangent space of a compact manifold.

To give a context to our results we need to explain in some detail the main ideas of the proof of the so-called stability conjecture: $C^{1}$ structurally stable diffeomorphisms are Axiom A and invariant submanifolds meet transversally, a result due to Mañé~\cite{kn:Mane}.

One of the main steps of the proof is that the $C^{1}$ structural stability implies that the closure of the set of periodic orbits is a hyperbolic invariant set. The key tool to prove this statement is the so-called Franks' lemma, we shall give an improved recent version of it for geodesic flows~\cite{kn:ARR}. Then, it is natural to expect that under this condition, the set of nonwandering points, the set where the dynamics is nontrivial, is exactly the closure of the set of periodic orbits. To show this the essential tool is the $C^{1}$ closing lemma proved by Pugh~\cite{kn:Pugh}, that is not available for geodesic flows up to date and is actually a very difficult problem in the theory of geodesic flows as we already mentioned in the Introduction (see~\cite{kn:Rifford}).

The step concerning the hyperbolicity of periodic orbits under stability assumptions has been extended and improved for geodesic flows in the context of the so-called \emph{Mañé perturbations}. Recall that a $C^{\infty}$ Hamiltonian $H:T^{*}M \longrightarrow \bbR$ defined in the cotangent bundle of $M$ is called \emph{Tonelli} if $H$ is strictly convex and superlinear in each tangent space $T_{\theta}T^{*}M$, $\theta \in T^{*}M$.

\begin{defi}\label{Mane-genericity}
A property $P$ of the Hamiltonian flow of a Tonelli Hamiltonian $H: T^{*}M \longrightarrow \bbR$ is called $C^{k}$ generic from Mañé's viewpoint if there exists a $C^{k}$ generic (from Baire's viewpoint) set of real valued functions $F_{P}$ defined on $M$ such that the Hamiltonian flow of $H_{f}(q,p) = H(q,p) + f(q)$ has the property $P$ for every $f \in F_{P}$. If the $C^{k}$-norm of $f$ is small, the Hamiltonian $H_{f}$ is called a $C^{k}$ Mañé perturbation of the Hamiltonian $H$.
\end{defi}

By the Maupertuis principle of classical mechanics. given a Riemannian metric $(M,g)$ defined in a compact manifold and a smooth scalar function $f : M \longrightarrow \bbR$ such that $\mid f(q) \mid <1$ for every $ q \in M$, the solutions of the Euler--Lagrange equation of the mechanical Lagrangian
\begin{align*}
L_{f}(q,p) &= \frac{1}{2}g_{q}(p.p) - f(q)
\\
\intertext{are the geodesics of the Riemannian metric}
g^{f}_{q}(p,p)&= 2(1- f(q))g_{q}(p,p).
\end{align*}

Therefore, by Legendre duality, every small $C^{k}$ Mañé perturbation $H_{f}=H(q,p) + f(q)$ of the Hamiltonian $H(q,p) = \frac{1}{2}g_{q}(p,p)$ defines the Riemannian Hamiltonian of a metric $g^{f}$ that is conformal to $g$ (see Arnold~\cite{kn:Arnold} for instance).

\begin{defi}\label{stability}
Given an integer $k\geq 2$, the geodesic flow $\phi_{t}$ of a compact Riemannian manifold $(M,g)$ is $C^{k}$ structural stable from Mañé's view point if there exists a $C^{k}$ small open neighborhood $V$ of the function $f_{0}(p) = 0$ for every $ p \in M$ in the set of real valued functions defined on $M$, such that the geodesic flow of the Hamiltonian $H_{f}(q,p) = \frac{1}{2}g_{q}(p,p) + f(p)$ is conjugate to $\phi_{t}$.
\end{defi}

This notion of structural stability from Mañé's viewpoint is stronger than the usual one, since it requires persistence dynamics in a neighborhood of special type of perturbations of the metric, not all perturbations of a metric are Mañé type perturbations. As we commented above, perturbations of $g$ which are not conformal to $g$ do not belong to the family of Mañé perturbations.

Applying techniques of control theory we obtain an extension of the hyperbolicity of the set of periodic orbits just considering Mañé perturbations.

\begin{theo}[{see~\cite{kn:ARR}}]\label{persistence}
Let $(M,g)$ be a compact manifold whose geodesic flow is $C^{2}$-structurally stable from Mañé's viewpoint. Then the closure of the set of periodic orbits is a hyperbolic set for the geodesic flow.
\end{theo}

As for the closing lemma, we have to rely on other kind of assumptions on the manifold to try to localize the set of periodic orbits. The natural domains to look for these assumptions are topology and global geometry.

\section{On the density of periodic orbits for rank one manifolds and manifolds with Gromov hyperbolic fundamental groups}\label{SECgromov}

Let us recall that a metric space $(X,d)$ is called \emph{geodesic} if every pair of points $xy$ can be joined by an isometric continuous embedding of an interval $c: [a,b] \longrightarrow X$ with $c(a)=x$, $c(b)=y$. The curve $c$ will be called a geodesic of the metric space, it corresponds to minimizing geodesics if $(X,d)$ is a Riemannian manifold. Let us denote by $[x,y]$ a minimizing geodesic joining $x$ to $y$ (there might be many).

\begin{defi}\label{Gromov-hyperbolic}
Given $\delta >0$ a complete geodesic metric space $(X,d)$ is called $\delta$-hyperbolic or Gromov hyperbolic if for every geodesic triangle $[x_{0},x_{1}] \cup [x_{1},x_{2}] \cup [x_{2},x_{0}]$ we have that the distance from every point $p \in [x_{i},x_{i+1}]$ to $[x_{i+1},x_{i+2}]\cup [x_{i+2},x_{i}]$ is bounded above by $\delta$ for every $i = 0,1,2$ (the indices are taken mod~3).
\end{defi}

\begin{defi}\label{visibility}
A complete Riemannian manifold $(N,g)$ without conjugate points is called a visibility manifold if given $p \in N$ and $\epsilon >0$ there exists $T>0$ such that for every pair of points $x, y $ in $N$, whenever the distance from $p$ to the geodesic $[x,y]$ is larger than $T$, then the angle subtended by the geodesics $[p,x]$ and $[p,y]$ at $p$ is less than $\epsilon$. If $T$ does not depend on $p$ the manifold $(N,g)$ is called a uniform visibility manifold.
\end{defi}

The universal covering of a compact surface without conjugate points and genus greater than one is a uniform visibility manifold, according to the work of Eberlein~\cite{kn:Eberlein}. In higher dimensions, the link between visibility manifolds and Gromov hyperbolic spaces is given in~\cite{kn:Ruggiero-Ensaios,kn:Ruggiero-SBM}.

\begin{theo}\label{visibility-gromov}
Let $(M,g)$ be a compact manifold without conjugate points. The universal covering is a uniform visibility manifold if and only if geodesic rays diverge and the fundamental group is Gromov hyperbolic.
\end{theo}

So the theory of Gromov hyperbolic spaces applies to visibility coverings of compact manifolds without conjugate points. Let us mention some of these results concerning the dynamics of the geodesic flow for the purposes of this article.

\begin{theo}[{see~\cite{kn:Eberlein}}]\label{rank-one-density}
Let $(M,g)$ be a compact manifold without conjugate points such that $(\tilde{M}, \tilde{g})$ is a visibility manifold. Then the geodesic flow is topologically transitive, and given any geodesic $\gamma \in \tilde{M}$ there exists a geodesic $\beta \in \tilde{M}$ that is bi-asymptotic to $\gamma$ having the following property:

There exists a sequence of axes $\beta_{n} \in \tilde{M}$, $n \rightarrow +\infty$, such that for every $t \in \bbR$,
\[
\lim_{n\rightarrow \infty}\beta_{n}(t) = \beta(t)
\]
where this limit is uniform on compact subsets. Whenever a geodesic $\beta$ is accumulated by axes in this way, we say that $\beta$ is in the closure $\bar{A}$ of the set of axes in the compact open topology of the continuous functions from $\bbR$ to $(\tilde{M}, \tilde{g})$, or simply, the compact open topology.
\end{theo}

\begin{theo}[{see~\cite{kn:Eberlein-96}}]\label{rank-one-nonpositive}

Let $(M,g)$ be a rank one compact manifold with nonpositive curvature. Then the set of periodic orbits of the geodesic flow is dense in $T_{1}M$.
\end{theo}

Combining Theorems~\ref{persistence} and~\ref{rank-one-nonpositive} we get

\begin{coro}\label{rank-one-stability}
Let $(M,g)$ be a compact manifold with nonpositive curvature. The geodesic flow is Anosov if and only if it is $C^{2}$-structurally stable from Mañé's viewpoint.
\end{coro}

\begin{proof}
The nontrivial part of the proof is the converse of the statement. If the geodesic flow is $C^{2}$-structurally stable from Mañé's viewpoint then by Theorem~\ref{persistence} the closure of the set of periodic orbits is hyperbolic. Since in manifolds with nonpositive curvature the intersections $H_{\theta}(0) \cap H_{-\theta}(0)$ generate flat, convex subsets of bi-asymptotic geodesics in $(\tilde{M}, \tilde{g})$, periodic geodesics have rank one. Therefore, Theorem~\ref{rank-one-nonpositive} yields that periodic orbits are dense and Theorem~\ref{persistence} then implies that the geodesic flow is Anosov.
\end{proof}

\section{The stability problem for surfaces and higher dimensional visibility manifolds}\label{sect.6}

The ``easy'' part of the proof of Theorem~\ref{main1} is the converse, namely, if the geodesic flow is Anosov then the geodesic flow is $C^{1}$-structurally stable (which is precisely Anosov's work). In particular, the geodesic flow is $C^{2}$ structurally stable from Mañé's viewpoint since the geodesic flows of small $C^{2}$ neighborhoods of conformal metrics are contained in $C^{1}$ small neighborhoods the geodesic flow of $(M,g)$.

\subsection{Structural stability for surfaces}
%\begin{proof}[{Proof of Theorem~\ref{main1}}]
Let us start the proof of the direct part of Theorem~\ref{main1} with the case of surfaces. The main steps of the proof in the two dimensional case are a sort of paradigm of what we would like to do in general dimensions. Let us summarize the main steps of the argument.
%\bigskip

%\textbf{Step 1}:
\begin{stepl}\label{step1}
The hyperbolicity of the closure of the set of periodic orbits determines the topology of the manifold: the surface has genus greater than one.
\end{stepl}
%\bigskip

To begin with, the compact surface $(M,g)$ cannot be a sphere because the sphere does not admit a Riemannian metric without conjugate points (otherwise its universal covering would be diffeomorphic to $\bbR^{2}$).

Next, assume that the geodesic flow of the compact surface $(M,g)$ without conjugate points is $C^{2}$ structurally stable from Mañé's viewpoint. By Theorem~\ref{persistence}, every closed orbit is a hyperbolic orbit of the geodesic flow. This implies that the surface has genus greater than one, since a torus without conjugate points is flat by the work of Hopf~\cite{kn:Hopf}.

Therefore, the surface $(M,g)$ has genus greater than one, its fundamental group is isomorphic to a Kleinian group and hence a Gromov hyperbolic group. This finishes the proof of Step~\ref{step1}.
%\bigskip

%\textbf{Step 2}:
\goodbreak
\begin{stepl}\label{step2}
The universal covering $(\tilde{M}, \tilde{g})$ is a visibility manifold.
\end{stepl}
%\bigskip

Once we know that the compact surface $(M,g)$ has genus greater than one we apply Eberlein's work~\cite{kn:Eberlein} as mentioned in the previous section.
%\bigskip

The next step is related to look at the obstructions to Anosov dynamics at the level of the topological dynamics of the geodesic flow of $(M,g)$. It is clear that the existence of bi-asymptotic geodesics in the universal covering is an obstruction to Anosov dynamics, since any Anosov geodesic flow is expansive and the only bi-asymptotic geodesic to a given geodesic $\gamma_{\theta} \subset \tilde{M}$ is $\gamma_{\theta}$ itself (see for instance~\cite{kn:Ruggiero-expansive}. So now we focus on the structure of the sets $S(\gamma_{\theta})$.
%\bigskip

%\textbf{Step 3}:
\begin{stepl}\label{step3}
Let $A$ be the union of the set of axes in $(\tilde{M}, \tilde{g})$, and let $\bar{A}$ be its closure in the compact open topology. Then, given a geodesic $\gamma_{\theta} \subset (\tilde{M}, \tilde{g})$ there exists $\gamma_{\bar{\theta}} \in \bar{A}$ such that
\[
S(\gamma_{\theta}) = S(\gamma_{\bar{\theta}}).
\]
%\bigskip

This follows from Theorem~\ref{rank-one-density}.
\end{stepl}
The last step of the proof follows from the hyperbolicity of the closure of the set of closed orbits and the connectedness of $S(\gamma_{\theta})$ (Corollary~\ref{connected-strip}).
%\bigskip

%\textbf{Step 4}:
\begin{stepl}\label{step4}
If the closure of the set of periodic orbits of the geodesic flow of $(M,g)$ is a hyperbolic set, then $S(\gamma_{\theta}) = \gamma_{\theta}$ for every $\theta \in T_{1}\tilde{M}$.
\end{stepl}
%\bigskip

Indeed, the proof of this statement is an extension of the proof of Lemma~\ref{hyperbolic-no-strip}. The assumption implies that the invariant set of orbits of the geodesic flow of $(\tilde{M}, \tilde{g})$ associated to the geodesics in $\bar{A}$ is a hyperbolic set. So let $\gamma_{\theta} \subset \tilde{M}$ be a geodesic having a bi-asymptotic geodesic $\gamma_{\bar{\theta}}$ in $\bar{A}$. By the work of Morse~\cite{kn:Morse}, the geodesics $\gamma_{\theta}$ and $\gamma_{\bar{\theta}}$ bound a strip in $\tilde{M}$ that is foliated by geodesics, all of them obviously bi-asymptotic to $\gamma_{\theta}$.

By the divergence of geodesic rays for surfaces and Lemma~\ref{asymptotic-1}, the intersection $H_{\theta}(0) \cap H_{-\theta}(0) = \Sigma(\gamma_{\theta})$ consists of a connected compact curve with boundary that contains $\gamma_{\bar{\theta}}(0)$ (up to an affine reparametrization of the latter geodesic).

By definition, this implies that we have a curve $I(\bar{\theta})$ homeomorphic to an interval in the intersection of $\tilde{W}^{s}(\bar{\theta})\cap \tilde{W}^{u}(\bar{\theta})$, this curve contains the points $\theta$ and $\bar{\theta}$ in $T_{1}\tilde{M}$. But we know by Lemma~\ref{asymptotic-1} that hyperbolicity implies that
\[
\tilde{W}^{s}(\bar{\theta})\cap \tilde{W}^{u}(\bar{\theta}) = \bar{\theta}
\]
since these invariant submanifolds intersect transversally at $\bar{\theta}$, and hence,
\begin{align*}
S(\gamma_{\bar{\theta}}) &= \gamma_{\bar{\theta}}.
\\
\intertext{Consequently,}
S(\gamma_{\theta}) &= S(\gamma_{\bar{\theta}}) = \gamma_{\bar{\theta}} = \gamma_{\theta}
\end{align*}
for every $\theta \in T_{1}\tilde{M}$.
%\bigskip

The proof of Theorem~\ref{main1} follows then from Step~\ref{step4} and Theorem~\ref{persistence}. The closure of the set of orbits subtending axes is a hyperbolic set, and the claim in Step~\ref{step4} tells us that the orbit of every $\theta $ in $T_{1}M$ itself is in the closure of the set of periodic orbits. Therefore, the whole set $T_{1}M$ is a hyperbolic set and the geodesic flow of $(M,g)$ is~Anosov.
%\end{proof}

\subsection{Structural stability for visibility manifolds}

From the proof of the structural stability for surfaces without conjugate points we notice the relevance of the Gromov hyperbolicity of the fundamental group and global geometry of the universal covering.
%%définir une entrée listée pour les steps référencés ci dessous
So let $(M,g)$ be a compact manifold without conjugate points such that $(\tilde{M}, \tilde{g})$ is a visibility manifold (actually a uniform visibility manifold by the compactness of $M$). Following the same line of reasoning for the surface case, observe that we do not need Step~\ref{step1} to prove Step~\ref{step2}, since it is an assumption in this case. Step~\ref{step3} follows from Theorem~\ref{visibility-gromov}, which ensures the quasi-convexity and the divergence of geodesic rays in $(M,g)$, Theorem~\ref{quasi-conv-continuity} to grant the uniform divergence of geodesic rays; and Theorem~\ref{rank-one-density}. The proof of Step~\ref{step4} mimics the proof for surfaces, replacing the work of Morse for surfaces by Theorems~\ref{visibility-gromov}, \ref{quasi-conv-continuity}, and Corollary~\ref{connected-strip} that ensures the connectedness of the intersection $H_{\theta}(0) \cap H_{-\theta}(0) =
\Sigma(\gamma_{\theta})$. Therefore, we have

\begin{theo}\label{stability-visibility}
Let $(M,g)$ be a compact manifold without conjugate points such that $(\tilde{M}, \tilde{g}))$ is a visibility manifold. The geodesic flow is Anosov if and only if it is $C^{2}$-structurally stable from Mañé's viewpoint.
\end{theo}


\section{The 3-dimensional case}\label{sect.7}

We shall show that the proof of Theorem~\ref{main1} in the three dimensional case follows the same line of reasoning of the proof of the two dimensional case. The roles of topology and hyperbolic global geometry of the manifold were crucial in the proof of Theorem~\ref{main1} item~\eqref{theo1.1.1}. This step of the proof for surfaces relies in their classification theory, that will be replaced by Thurston's geometrization theory in the case of three dimensional manifolds.

\subsection{Hyperbolic periodic orbits and the fundamental group}

Let us start with the following result linking hyperbolic closed geodesics and a special algebraic property of Gromov hyperbolic groups: the absence of higher rank abelian subgroups:

\begin{lemm}\label{Preissmann}
Let $(M,g)$ be a compact manifold without conjugate points such that $(\tilde{M}, \tilde{g})$ is quasi-convex and geodesic rays diverge. Then, if every closed orbit is hyperbolic the fundamental group is a Preissmann group, namely, every nontrivial abelian subgroup is infinite cyclic.
\end{lemm}

\begin{proof}
The proof relies on Lemma~\ref{hyperbolic-no-strip} and basic properties of the fundamental group of manifolds without conjugate points. Let us remind such properties briefly.

Let $\Gamma \subset \pi_{1}(M)$ be an abelian subgroup of the fundamental group of $M$. The group $\pi_{1}(M)$ has no torsion and hence $\Gamma$ is isomorphic to the product of a finite number of subgroups isomorphic to $\bbZ$.

Let $h_{1}, h_{2}$ be two nontrivial elements in $\Gamma$, and let $\gamma_{\theta_{1}}$ be an axis of $h_{1}$, so
\[
h_{1}(\gamma_{\theta_1}(t)) = \gamma_{\theta_1}(t + T_{1})
\]
for every $t \in \bbR$, where $T_{1}$ is a period of $\pi(\gamma_{\theta_{1}})$.
%\bigskip

\begin{enonce*}[plain]{Claim}
The geodesic $h_{2}(\gamma_{1}(t))$ is also an axis of $h_{1}$.
\end{enonce*}
%\bigskip

This is a simple consequence of the fact that $h_{1}\circ h_{2} = h_{2} \circ h_{1}$:
\[
h_{1}(h_{2}(\gamma_{\theta_{1}}(t))) = h_{2}(h_{1}(\gamma_{\theta_{1}}(t))) = h_{2}(\gamma_{\theta_{1}}(t+ T_{1}))
\]

so the geodesic $h_{2}(\gamma_{\theta_{1}}(t))$ is preserved by $h_{1}$ as well. This proves the claim.
%\bigskip

Now, Corollary~\ref{connected-strip} and Lemma~\ref{hyperbolic-no-strip} imply that the geodesics $\gamma_{\theta_{1}}$ and $h_{2}(\gamma_{\theta_{1}}) $ must coincide (as sets). Since $\pi_{1}(M)$ has no torsion, we deduce that $h_{1}$ and $h_{2}$ are powers of a single element $h_{0}$ in $\pi_{1}(M)$, preserving the geodesic $\gamma_{\theta_{1}}$. This clearly yields that the group $\Gamma$ has to be infinite cyclic.
\end{proof}

The Preissmann property is satisfied by Gromov hyperbolic groups, but it is not enough to characterize such groups.

\subsection{Three dimensional manifolds without conjugate points and Thurston's geometrization}

A challenging question is whether the Preissmann property is sufficient to characterize the fundamental group of compact manifolds without conjugate points. The next result proved in~\cite[Chapter~8]{kn:Ruggiero-Ensaios}, is a partial positive answer to this question.

\begin{theo}\label{Preissmann-3}
Let $(M,g)$ be a compact Riemannian manifold without conjugate points with dimension~3. Suppose that the fundamental group satisfies the Preissmann property. Then $\pi_{1}(M)$ is a Gromov hyperbolic group and $M$ admits a metric of constant negative curvature.
\end{theo}

Now we can prove Theorem~\ref{main1} item~\eqref{theo1.1.2}:
%\bigskip

Let $(M,g)$ satisfy the assumptions of Theorem~\ref{main1}. $(\tilde{M}, \tilde{g})$ is quasi-convex and geodesic rays diverge. Moreover, the geodesic flow is $C^{2}$ structurally stable from Mañé´s viewpoint. By Theorem~4.7 the closure of the set of periodic orbits is a hyperbolic set for the geodesic flow. In particular, every periodic orbit is hyperbolic and by Lemma~\ref{Preissmann} the fundamental group is a Preissmann group. Therefore, by Theorem~\ref{Preissmann-3} the fundamental group is Gromov hyperbolic, and hence Theorem~\ref{visibility-gromov} implies that the universal covering is a visibility manifold. Now, we deduce that the geodesic flow of $(M,g)$ is Anosov by Theorem~\ref{stability-visibility}.

\section{The stability problem in higher dimensions}\label{SEChigher}

The study of the $C^{2}$ structural stability from Mañé's viewpoint in higher dimensions without an appropriate closing lemma and no clue about the nature of the fundamental group requires a different strategy. Theorem~\ref{main2} will follow from the extension of Theorem~\ref{rank-one-nonpositive} to our context.

\begin{theo}\label{density2}
Under the assumptions of Theorem~\ref{main2} the set of periodic orbits is dense in $T_{1}M$.
\end{theo}

Throughout the section, $(M,g)$ will be a compact Riemannian manifold without conjugate points whose universal covering is $K,C$-quasi-convex where geodesic rays diverge, having a dense set of generalized rank one points.

We know that the set of recurrent orbits has total Liouville measure in $T_{1}M$. Since the set of points of $T_{1}M$ where $W^{s}_{loc}(\theta) \cap W^{u}_{loc}(\theta) = \{ \theta\}$ contains an open and dense set by the assumption of Theorem~\ref{main2}, the set of recurrent points with this property is dense as well.

So let $\eta$ be such a point, let $B(\eta) \subset T_{1}M$ be an open ball around $\eta$ of generalized rank one points.

\subsection{Local product structure for the geodesic flow in neighborhoods of generalized rank one points}

The definition of the generalized rank one assumption implies the existence of a special cross section $\Sigma_{\eta}$ for the geodesic flow containing $\eta$ foliated by local unstable sets. Namely, given $\epsilon >0$ and $\theta =(p,v) \in T_{1}\tilde{M}$, let $\Psi^{\theta} : D_{\epsilon} \longrightarrow B_{\epsilon}^{h}(\theta)$ be the exponential map of the horosphere $H_{\theta}(0)$, where $D_{\epsilon}$ is the open disk of radius $\epsilon$ around $0 \in T_{p}H_{\theta}(0)$, and $B_{\epsilon}^{h}(\theta)$ is the geodesic ball of radius $\epsilon$ in $H_{\theta}(0)$ centered at $p$, with respect to the restriction of $\tilde{g}$ to $H_{\theta}(0)$. By the continuity of $H_{\theta}(0)$ with respect to $\theta$ (Theorem~\ref{divergence-continuity}), we can choose $\epsilon >0$ such that $\Psi^{\theta}$ is a diffeomorphism for every $\theta$ in a fundamental domain of $T_{1}\tilde{M}$. Let us define
\[
\tilde{W}^{s}_{\epsilon}(\theta) = \left\{\left(x,-\nabla_{x}b^{\theta}\right),\, x \in B_{\epsilon}^{h}(\theta)\right\},
\]
\[
\tilde{W}^{u}_{\epsilon}(\theta) = \left\{\left(x, \nabla_{x}b^{-\theta}\right),\,x \in B_{\epsilon}^{h}(-\theta)\right\}.
\]
The above sets are homeomorphic to $(m-1)$ dimensional open balls for every $\theta \in T_{1}\tilde{M}.$

There exists $\epsilon_{0}>0$ such that the projections
\begin{align*}
W^{s}_{\epsilon_{0}}(\theta) &= \Pi\left(\tilde{W}^{s}_{\epsilon_{0}}(\theta)\right),
\\
W^{u}_{\epsilon_{0}}(\theta) &= \Pi\left(\tilde{W}^{u}_{\epsilon_{0}}(\theta)\right)
\end{align*}
are embedded continuous submanifolds in $T_{1}M$. Reducing $\epsilon_{0}$ if necessary, we can suppose that $W^{s}_{\epsilon_{0}}(\eta) \subset B(\eta)$. To simplify notation we just identify $\epsilon_{0}$ with $\epsilon$. Since every point in $B(\eta)$ is a generalized rank one point, the set
\[
\Sigma_{\eta}^{us} = \bigcup_{\theta\,\in\,W^{s}_{\epsilon}(\eta)} W^{u}_{\epsilon}(\theta)
\]
can be continuously parametrized as the homeomorphic image of the product $D^{m-1}_{\epsilon} \times D^{m-1}_{\epsilon}$, where $D^{m-1}_{\epsilon}$ is the open disk of radius $\epsilon$ in $\bbR^{m-1}$, $m$ being the dimension of $M$.

The set
\[
U_{\epsilon}^{us}(\eta) = \bigcup_{\theta\,\in\,W^{s}_{\epsilon}(\eta),\,\mid t \mid\,<\,\epsilon} \phi_{t}\left(W^{u}_{\epsilon}(\theta)\right)
\]
that contains $\Sigma_{\eta}^{us}$, can be parametrized as the homeomorphic image of $D^{m-1}_{\epsilon} \times D^{m-1}_{\epsilon}\times\{ \mid t \mid <\epsilon\} $. So by Brouwer's open mapping theorem it is a $(2m-1)$-dimensional open subset of $T_{1}M$ and $\Sigma_{\eta}^{us}$ is a continuous cross section for the geodesic flow (see also~\cite{kn:Ruggiero-expansive} for a very similar construction).

The same happens with the section
\[
\Sigma_{\eta}^{su} = \bigcup_{\theta\,\in\,W^{u}_{\epsilon}(\eta)} W^{s}_{\epsilon}(\theta)
\]
and the set
\[
U_{\epsilon}^{su}(\eta) = \bigcup_{\theta\,\in\,W^{u}_{\epsilon}(\eta),\,\mid t \mid\,<\,\epsilon} \phi_{t}\left(W^{s}_{\epsilon}(\theta)\right)
\]
that is an open neighborhood of $\eta$ by the same above argument.


\begin{lemm}\label{product-structure}
The open neighborhood of $\eta$ given by
\[
U_{\epsilon}(\eta)= U_{\epsilon}^{us}(\eta) \cap U_{\epsilon}^{su}(\eta)
\]
has a local product structure in the following sense: for each point $\theta \in U_{\epsilon}(\eta)$ there exist $\theta_{s} \in W^{s}_{\epsilon}(\eta)$, $\theta_{u} \in W^{u}_{\epsilon}(\eta)$, continuous real valued functions $b_{s}(\theta), b_{u}(\theta)$ close to zero such that
\[
\theta = W^{s}_{loc}\left(\phi_{b_{u}}(\theta_{u}\right) \cap W^{u}_{loc}\left(\phi_{b_{s}}(\theta_{s})\right).
\]
\end{lemm}

\begin{proof}
The statement is more or less self evident from the definitions of the sets. Both $U_{\epsilon}^{us}(\eta), U_{\epsilon}^{su}(\eta)$ are open neighborhoods of $ W^{s}_{\epsilon}(\eta) \cup W^{u}_{\epsilon}(\eta)$ so their intersection is obviously a nonempty open set. The statement follows then from the very definition of these neighborhoods.
\end{proof}


\subsection{Local expansiveness of the geodesic flow near generalized rank one points}

Since $\eta$ is recurrent (meaning forward and backward recurrent) there exists a sequence $t_{n} \rightarrow +\infty$, with $t_{0}=0$, such that $\phi_{t_{n}}(\eta) \in \Sigma_{\eta}^{us}$ and converges to $\eta$. By the smooth dependence of solutions of ordinary differential equations with respect to initial conditions, there exist an open subset $\Sigma_{\eta}^{n}$ of $\Sigma_{\eta}^{us}$ (in the relative topology) and a smooth function $\bar{t}_{n} : \Sigma_{\eta}^{n} \longrightarrow \bbR$ satisfying $\bar{t}_{n}(\eta) = t_{n}$, such that the orbit of every $\theta \in \Sigma_{\eta}^{n}$ hits $\Sigma_{\eta}^{us}$ at $\phi_{\bar{t}_{n}(\theta)}(\theta)$.

So let $P_{n} : \Sigma_{\eta}^{n} \longrightarrow \Sigma_{\eta}^{us}$ be the return map given by $P_{n}(\theta) =\phi_{\bar{t}_{n}(\theta)}(\theta) $.

Let $W^{s}_{loc}(\theta)$, $W^{u}_{loc}(\theta)$, be the connected components of the intersections of $W^{s}(\theta)$, $W^{u}(\theta)$ with $U_{\epsilon}^{us}(\eta)$ which contain $\theta \in U_{\epsilon}^{us}(\eta)$.


By the above definitions we have that
\begin{align*}
P_{n}\left(W^{u}_{\epsilon}(\theta) \cap \Sigma_{\eta}^{n}\right)&\subset W^{u}_{loc}\left(\phi_{\bar{t}_{n}(\theta)}(\theta)\right)
\\
\intertext{for every $\theta \in \Sigma^{n}_{\eta}$ and}
P_{n}\left(W^{s}_{\epsilon}(\eta) \cap \Sigma_{\eta}^{n}\right)&\subset W^{s}_{loc}\left(\phi_{t_{n}}(\eta)\right),
\end{align*}
for every $n \geq 0$.
%\bigskip
\pagebreak
\begin{lemm}\label{contraction}
Under the assumptions of Theorem~\ref{main2}, let $\eta$, $P_{n}$, $W^{s}_{\epsilon}(\eta)$, be defined as above. Given $0 <\delta <\epsilon$, there exists $n_{\delta}>0$ such that for every $n >n_{\delta}$
\begin{enumerate}
\item \label{lemm8.3.1}$P_{n}$ is defined in $W^{s}_{\epsilon}(\eta)$,
\item \label{lemm8.3.2} $P_{n}(W^{s}_{\epsilon}(\eta))$ is strictly contained in $W^{s}_{\delta}(\phi_{t_{n}}(\eta)).$
\end{enumerate}
\end{lemm}
\begin{proof}
Indeed, otherwise there exists $\delta_{0} < \epsilon$ such that for each $k>0$ there exist $n_{k} \geq k$ and a point $\tau_{k} \in W^{s}_{\epsilon}(\eta)$ such that $P_{n_{k}}(\tau_{k}) $ is either not defined or defined but outside $W^{s}_{\delta_{0}}(\phi_{t_{n_{k}}}(\eta))$. In any case the orbit of $\tau_{k}$ satisfies
\[
\phi_{t_{n_{k}}}(\tau_{k}) \in W^{s}\left(\phi_{t_{n_{k}}}(\eta)\right) - W^{s}_{\delta_{0}}\left(\phi_{t_{n_{k}}}(\eta)\right).
\]

Let us consider a lift $\gamma_{\bar{\eta}}$ of $\gamma_{\eta}$ in $\tilde{M}$, with $\bar{\eta}= (p,v)$, and a lift $\gamma_{\bar{\tau}_{k}}$ of $\gamma_{\tau_{k}}$ in $\tilde{M}$. Since $\tau_{k} \in W^{s}_{\epsilon}(\eta)$, the lift $\gamma_{\bar{\tau}_{k}}$ can be chosen asymptotic to $\gamma_{\bar{\eta}}$. Moreover, if $\tau_{k} = (p_{k}, v_{k})$ then $p_{k} \in H_{\bar{\eta}}(0)$.

By Lemma~\ref{asymptotic-quasi-convex}
\[
d\left(\gamma_{\bar{\eta}}(t), \gamma_{\bar{\tau}_{k}}(t)\right) \leq Kd(p_{k}, p) +C
\]
for every $t \geq 0$,
and the contradiction assumption yields that there exists $a = a(\delta_{0}) >0$ such that
\[
d\left(\gamma_{\bar{\tau}_{k}}(t_{n_{k}}), \gamma_{\bar{\eta}}(t_{n_{k}})\right) \geq a
\]
for every $k>0 $.

Since $\eta$ is recurrent, there exists a sequence of covering isometries $T_{k}: \tilde{M} \longrightarrow \tilde{M}$ such that
\begin{enumerate}
\item \label{proof-lemm8.3.1} The pairs $(T_{k}(\gamma_{\bar{\eta}}(t_{n_{k}})), dT_{k}(\gamma_{\bar{\eta}}'(t_{n_{k}})))$ converge to $\bar{\eta}$.
\item \label{proof-lemm8.3.2} The sequence $T_{k}(\gamma_{\bar{\tau}_{k}}(t_{n_{k}})) $ is contained in a compact ball centered at $p$,
\item \label{proof-lemm8.3.3} A convergent subsequence $n_{k_{j}}$ of the pairs $(T_{k}(\gamma_{\bar{\eta}}(t_{n_{k}})), dT_{k}(\gamma_{\bar{\eta}}'(t_{n_{k}})))$ gives rise to a geodesic $\beta$ defined by
\[
\beta(t) = \lim_{j\,\rightarrow\,\infty} T_{k_{j}}\left(\gamma_{\bar{\tau}_{k_{j}}}\left(t+t_{n_{k_{j}}}\right)\right)
\]
that is bi-asymptotic to $\gamma_{\bar{\eta}}$.
\item $d(\beta(0), \gamma_{\bar{\eta}}(0)) \geq a >0$.
\end{enumerate}

Thus, we get a geodesic $\beta$ different from $\gamma_{\bar{\eta}}$ that is bi-asymptotic to $\gamma_{\bar{\eta}}$. But by Corollary~\ref{connected-strip} this would generate a connected subset in the intersection $W^{s}_{loc}(\eta) \cap W^{u}_{loc}(\eta)$ containing more than one point, which is impossible by the generalized rank one definition and the choice of $\eta$. This finishes the proof of the Claim.
\end{proof}
%\bigskip

A statement analogous to Lemma~\ref{contraction} proceeds for $W^{u}_{\epsilon}(\eta)$ and $s_{n} \rightarrow -\infty$ such that $\phi_{s_{n}}(\eta) \rightarrow \eta$, by the reversibility of the geodesic flow.

\begin{lemm}\label{contraction2}
Under the assumptions of Theorem~\ref{main2}, let $\eta$, $s_{n}$ be as above, $P_{n}^{-}$ the return map of $\Sigma_{\eta}^{su}$ relative to $s_{n}$. Given $0 <\delta <\epsilon$, there exists $n_{\delta}^{-}<0$ such that for every $n <n_{\delta}^{-}$
\begin{enumerate}
\item \label{lemm8.4.1}$P_{n}^{-}$ is defined in $W^{u}_{\epsilon}(\eta)$,
\item \label{lemm8.4.2}$P_{n}^{-}(W^{u}_{\epsilon}(\eta))$ is strictly contained in $W^{u}_{\delta}(\phi_{s_{n}}(\eta)).$
\end{enumerate}
\end{lemm}


\subsection{Proof of Theorem~\ref{density2}}

The set $W^{s}_{\epsilon}(\eta)$ is a subset of the open neighborhood $B(\eta)$ where every point is generalized rank one. As observed in Subsection~8.1, each intersection of the type $W^{u}_{\epsilon} (\theta) \cap W^{s}_{\epsilon}(\eta)$ consists of a single point for every $\theta \in \Sigma_{\eta}^{us}$. This allows us to define a continuous projection
\begin{align*}
&\Pi_{s} : \Sigma_{\eta} \longrightarrow W^{s}_{\epsilon}(\eta)
\\
&\Pi_{s}(\theta) = W^{u}_{loc} (\theta) \cap W^{s}_{\epsilon}(\eta)
\end{align*}
along the foliation of $\Sigma_{\eta}$ formed by the sets $W^{u}_{loc}(\theta)$.


By Lemma~\ref{contraction}, the fact that $\lim_{n\rightarrow \infty}\phi_{t_{n}}(\eta) \rightarrow \eta$, and the continuity of the sets $W^{s}_{r}(\sigma)$ with respect to $\sigma$ and $r$, we have that given $0 < \delta <\epsilon$, there exists $ n(\delta)>0$ such that
\[
\Pi_{s}\left(P_{n}\left(W^{s}_{\epsilon}(\eta)\right)\right) \subset W^{s}_{\delta}(\eta)
\]
for every $ n \geq n(\delta)$. Therefore, the map $\Pi_{s} \circ P_{n} : W^{s}_{\epsilon}(\eta) \longrightarrow W^{s}_{\epsilon}(\eta)$ sends a set homeomorphic to a $(m-1)$ dimensional ball into itself.

By Brouwer's fixed point theorem, there exists a fixed point $\sigma_{s,\,\delta} \in W^{s}_{\epsilon}(\eta)$ of $\Pi_{s} \circ P_{n}$.

An analogous argument applying Lemma~\ref{contraction2} shows that there exists a fixed point $\sigma_{u,\,\delta}$ for $P_{n_{\delta}^{-}}^{-}\circ \Pi_{u}$ in $W^{u}_{\epsilon}(\eta)$, where $\Pi_{u} : \Sigma_{\eta}^{su} \longrightarrow W^{u}_{\epsilon}(\eta)$ is the projection along stable leaves onto $W^{u}_{\epsilon}(\eta)$.
%\bigskip

\begin{enonce*}[plain]{Claim~1} If $\delta$ is small enough, there exists $\sigma_{\delta} \in U_{\epsilon}(\eta)$ such that
\[
\sigma_{\delta} = W^{s}_{\epsilon}\left(\phi_{r(u)}(\sigma_{u, \,\delta})\right)\cap W^{u}\left(\phi_{r(s)}(\sigma_{s,\,delta})\right),
\]
for some small numbers $r(s), r(u)$ depending on $\sigma$ and $\delta$.
\end{enonce*}
%\bigskip

Indeed, if $\delta $ is small then $W^{s}_{\epsilon}(\sigma_{u, \,\delta})$ is very close to $W^{s}_{\epsilon}(\eta)$ by the continuity of invariant foliations (Theorem~\ref{divergence-continuity}). So the continuous submanifolds
\begin{align*}
W^{cs}_{\epsilon}(\eta) &= \bigcup_{\left\{\mid t \mid\,<\,\epsilon\right\}} \phi_{t}\left(W^{s}_{\epsilon}(\eta)\right),
\\
W^{cs}_{\epsilon}(\sigma_{u,\,\delta}) &= \bigcup_{\left\{\mid t \mid\,<\,\epsilon\right\}} \phi_{t}\left(W^{s}_{\epsilon}(\sigma_{u,\,\delta})\right),
\end{align*}
will be close to each other as well and hence, the projection
\[
\Pi_{s} : W^{cs}_{\epsilon}(\sigma_{u,\,\delta}) \cap U_{\epsilon}(\eta) \longrightarrow W^{cs}_{\epsilon}(\eta)
\]
along the unstable sets is a homeomorphism onto its image that has to be close to $W^{cs}_{\epsilon}(\eta) \cap U_{\epsilon}(\eta)$ by the continuity of unstable sets and the small size of $\delta$ (if $\delta =0$ the projection is just the identity map).

Therefore, the set $\Pi_{s}(W^{cs}_{\epsilon}(\sigma_{u,\,\delta})\cap U_{\epsilon}(\eta)) $ covers almost all $W^{cs}_{\epsilon}(\eta) \cap U_{\epsilon}(\eta) $, so reducing $\delta$ if necessary this set contains $\sigma_{s,\,\delta}$ and will be hit by the unstable set of some point $\sigma_{\delta}$ in the orbit of $W^{s}_{\epsilon}(\sigma_{u,\,\delta})$ as claimed.
%\bigskip

\begin{enonce*}[plain]{Claim~2}: The orbit of $\sigma_{\delta}$ is periodic.
\end{enonce*}
%\bigskip

Indeed, the point $\sigma_{\delta}$ has the same unstable local set of $\phi_{r(s)}(\sigma_{u,\,\delta})$, whose unstable set is periodic by $\phi_{t_{n_{\delta}}+ r(s)}$ because the unstable set of $\sigma_{u,\,\delta}$ is periodic by $\phi_{t_{n_{\delta}}}$. At the same time, $\sigma_{\delta}$ has the same stable local set of $\phi_{r(u)}(\sigma_{s,\,\delta})$, and the stable set of this point is periodic by $\phi_{s_{n_{\delta}^{-}}+ r(u)}$. Since the geodesic flow preserves invariant sets, the point $\sigma_{\delta}$ must be periodic with period $t_{n_{\delta}}+ r(s) + s_{n_{\delta}^{-}} + r(u)$.
%\bigskip

To finish the proof of Theorem~\ref{density2}, observe that we can approach $\eta$ as close as desired by making $\delta \rightarrow 0$, and since recurrent points are dense, we conclude that periodic points are also dense.
%\bigskip

\begin{proof}[{Proof of Theorem~1.2}]
%\bigskip
Theorem~\ref{main2} follows from Theorem~\ref{density2} and the fact that the closure of the set of periodic orbits is hyperbolic by Theorem~\ref{persistence}.
\end{proof}


\begin{thebibliography}{99}
\bibitem{kn:Anosov} [Ano69] Anosov, Dmitriĭ V. Geodesic flows on closed Riemannian manifolds of negative curvature, Proceedings of the Steklov Institute of Mathematics, 90, American Mathematical Society, 1969 (Translated from the Russian by S. Feder)
\bibitem{kn:Arnold} [Arn78] Arnold, Vladimir I. Mathematical Methods of Classical Mechanics, Graduate Texts in Mathematics, 60, Springer, 1978
\bibitem{kn:Ballmann} [Bal95] Ballmann, Werner Lectures on spaces of nonpositive curvature, DMW Seminar, 25, Birkhäuser, 1995
\bibitem{kn:BBB} [BBB87] Ballmann, Werner; Brin, Michael; Burns, Keith On surfaces with no conjugate points, J. Differ. Geom., Volume 25 (1987), pp. 249-273
\bibitem{kn:BGS} [BGS85] Ballmann, Werner; Gromov, Michael L.; Schroeder, Viktor Manifolds of nonpositive curvature, Progress in Mathematics, 61, Birkhäuser, 1985
\bibitem{kn:Burns} [Bur92] Burns, Keith The flat strip theorem fails for surfaces with no conjugate points, Proc. Am. Math. Soc., Volume 115 (1992), pp. 199-206
\bibitem{kn:CS} [CS86] Croke, Christopher B.; Schroeder, Viktor The fundamental group of compact manifolds without conjugate points, Comment. Math. Helv., Volume 61 (1986) no. 1, pp. 161-175
\bibitem{kn:Eberlein} [Ebe72] Eberlein, Patrick B. Geodesic flows in certain manifolds without conjugate points, Trans. Am. Math. Soc., Volume 167 (1972), pp. 151-170
\bibitem{kn:Eberlein-96} [Ebe96] Eberlein, Patrick B. Geometry of nonpositively curved manifolds, Chicago Lectures in Mathematics, University of Chicago Press, 1996
\bibitem{kn:Eschenburg} [Esc77] Eschenburg, Jost-Hinrich Horospheres and the stable part of the geodesic flow, Math. Z., Volume 153 (1977), pp. 237-251
\bibitem{kn:Green} [Gre54] Green, Leon W. Surfaces without conjugate points, Trans. Am. Math. Soc., Volume 76 (1954), pp. 529-546
\bibitem{kn:Gromov} [Gro87] Gromov, Mikhael L. Hyperbolic Groups, Essays in Group Theory (Mathematical Sciences Research Institute Publications), Volume 8, Springer, 1987, pp. 75-264
\bibitem{kn:Hopf} [Hop48] Hopf, Eberhard Closed surfaces without conjugate points, Proc. Natl. Acad. Sci. USA, Volume 34 (1948), pp. 47-51
\bibitem{kn:HPS} [HPS77] Hirsch, Morris; Pugh, Charles C.; Shub, Michael Invariant Manifolds, Lecture Notes in Mathematics, 583, Springer, 1977
\bibitem{kn:Klingenberg} [Kli74] Klingenberg, Wilhelm Riemannian manifolds with geodesic flows of Anosov type, Ann. Math., Volume 99 (1974), pp. 1-13
\bibitem{kn:ARR} [LRR16] Lazrag, Ayadi; Rifford, Ludovic; Ruggiero, Rafael Franks’ Lemma for $${C}^{2}$$ Mañé perturbations of Riemannian metrics and applications to persistence, J. Mod. Dyn., Volume 10 (2016), pp. 379-411
\bibitem{kn:Mane} [Mañ82] Mañé, Ricardo An ergodic closing lemma, Ann. Math., Volume 116 (3) (1982), pp. 503-540
\bibitem{kn:Mane2} [Mañ87] Mañé, Ricardo On a theorem of Klingenberg, Dynamical systems and bifurcation theory (Rio de Janeiro, 1985) (M. Camacho, F. Takens Editors M. Pacífico, ed.) (Pitman Research Notes in Mathematics Series), Volume 160, Longman Scientific \& Technical, Harlow (1987), pp. 319-345
\bibitem{kn:Mane88} [Mañ88] Mañé, Ricardo A proof of the $${C}^1$$ stability conjecture, Publ. Math., Inst. Hautes Étud. Sci., Volume 66 (1988), pp. 161-210
\bibitem{kn:Morse} [Mor24] Morse, Harold Marston A fundamental class of geodesics on any closed surface of genus greater than one, Trans. Am. Math. Soc., Volume 26 (1924), pp. 25-60
\bibitem{kn:Newhouse} [New77] Newhouse, Sheldon E. Quasi-elliptic periodic points in conservative dynamical systems, Am. J. Math., Volume 99 (1977) no. 5, pp. 1061-1087
\bibitem{kn:Pesin} [Pes77] Pesin, Ja B. Geodesic flows on closed Riemannian manifolds without focal points, Math. USSR, Izv., Volume 11 (1977) no. 6, pp. 1195-1228
\bibitem{kn:Pugh} [PR83] Pugh, C. Charles; Robinson, Clark The $${C}^{1}$$ closing lemma, including Hamiltonians, Ergodic Theory Dyn. Syst., Volume 3 (1983), pp. 261-313
\bibitem{kn:Rifford} [Rif12] Rifford, Ludovic Closing geodesics in $${C}^{1}$$ topology, J. Differ. Geom., Volume 91 (2012) no. 3, pp. 361-382
\bibitem{kn:Robin-Robinson} [Rob70a] Robinson, Clark Generic properties of conservative systems, Am. J. Math., Volume 92 (1970), pp. 562-603
\bibitem{kn:Robin-Robinson2} [Rob70b] Robinson, Clark Generic properties of conservative systems. II., Am. J. Math., Volume 92 (1970), pp. 897-906
\bibitem{kn:Ruggiero-SBM} [Rug94] Ruggiero, Rafael Expansive dynamics and hyperbolic geometry, Bol. Soc. Bras. Mat., Nova Sér., Volume 25 (1994) no. 2, pp. 139-172
\bibitem{kn:Ruggiero-expansive} [Rug97] Ruggiero, Rafael Expansive geodesic flows in manifolds without conjugate points, Ergodic Theory Dyn. Syst., Volume 17 (1997), pp. 211-225
\bibitem{kn:Ruggiero-Asterisque} [Rug03] Ruggiero, Rafael On the divergence of geodesic rays in manifolds without conjugate points, dynamics of the geodesic flow and global geometry, Geometric methods in dynamics (II). Volume in honor of Jacob Palis. In part papers presented at the international conference on dynamical systems held at IMPA, Rio de Janeiro, Brazil, July 2000, to celebrate Jacob Palis’ 60th birthday (Astérisque), Volume 287, Société Mathématique de France, 2003, pp. 231-249
\bibitem{kn:Ruggiero-Ensaios} [Rug07] Ruggiero, Rafael Dynamics and global geometry of manifolds without conjugate points, Ensaios Matemáticos, 12, Sociedade Brasileira de Matemática, 2007
\bibitem{kn:Ruggiero-div} [Rug08] Ruggiero, Rafael A note on the divergence of geodesic rays in manifolds without conjugate points, Geom. Dedicata, Volume 134 (2008), pp. 131-138
\end{thebibliography}
\end{document}
